\documentclass[11pt]{article}
\usepackage{amsmath,amssymb,amsthm}

\title{Picard Groups and Cohomology on $\mathrm{Spec}(R)$}
\author{UlamLib seed note}
\date{}

\newtheorem{definition}{Definition}[section]
\newtheorem{lemma}[definition]{Lemma}
\newtheorem{theorem}[definition]{Theorem}
\newtheorem{corollary}[definition]{Corollary}

\begin{document}
\maketitle

\section{Scope}

This follow-on note packages the cohomological refinement of the Picard track.
The intended Lean destination is
\texttt{UlamLib/AlgebraicGeometry/PicardCohomology.lean}.

\paragraph{Target Lean declarations.}
This note should eventually feed declarations with names close to
\texttt{picardToH1}, \texttt{invertibleSheafToUnitsCocycle}, and
\texttt{picardSpecH1Equiv}.

\section{Ambient assumptions}

Fix a commutative ring $R$ and work on $\mathrm{Spec}(R)$ with the standard
structure sheaf and its unit sheaf.

\section{Declaration blueprint}

\begin{enumerate}
\item \texttt{invertibleSheafToUnitsCocycle}: build a cocycle from an
invertible sheaf.
\item \texttt{unitsCocycleToInvertibleSheaf}: the reverse construction.
\item \texttt{picardToH1}: map from Picard classes to sheaf cohomology.
\item \texttt{picardSpecH1Equiv}: the final identification.
\end{enumerate}

\section{Seed theorem cluster}

\begin{lemma}[Tensor product respects cocycles]\label{lem:tensor-respects-cocycles}
The cocycle attached to a tensor product of invertible sheaves is the product
of the corresponding unit-valued cocycles.
\end{lemma}

\begin{theorem}[Picard to cohomology]\label{thm:picard-to-cohomology}
There is a natural map from the Picard group of $R$ to
$H^1(\mathrm{Spec}(R), \mathcal{O}_X^\times)$.
\end{theorem}

\begin{theorem}[Cohomological identification]
\label{thm:cohomological-identification}
The natural map from the Picard group of $R$ to
$H^1(\mathrm{Spec}(R), \mathcal{O}_X^\times)$ is an isomorphism.
\end{theorem}

\section{Formalization order}

\begin{enumerate}
\item package the unit-valued cocycle constructions
\item prove compatibility with tensor product
\item define the Picard-to-cohomology map
\item prove the final equivalence
\end{enumerate}

\section{Suggested first PR}

A good first contribution is to split one of the large statements into a helper
lemma plus a theorem, for example separating tensor compatibility from the
construction of the Picard-to-cohomology map.

\end{document}
